\documentclass[letterpaper]{article}
    \usepackage[top=.6in,bottom=.6in,left=1.1in,right=1.1in]{geometry}
    \usepackage{hyperref} % pdf metadata
    \usepackage{tagging}
    \usepackage{array}    % used by tables
    \usepackage{fontspec}
    \usepackage{ulem}  % underlining
    \usepackage{graphicx}
    \usepackage{eso-pic}

    \newfontfamily\Lato[
        Path = fonts/,
        Extension = .ttf,
        UprightFont = *-Regular,
        BoldFont = *-Bold,
        ItalicFont = *-Italic,
        BoldItalicFont = *-BoldItalic,
        Ligatures = NoCommon,
]{Lato}

    \pagestyle{empty}
    \setmainfont{Lato}
    \raggedright

    % Ragged right in tables
    %   see: https://tex.stackexchange.com/questions/12703/how-to-create-fixed-width-table-columns-with-text-raggedright-centered-raggedlef
    \newcolumntype{L}[1]{>{\raggedright\let\newline\\\arraybackslash}p{#1}}



%%%%%%%%%%%%%%%%%%%%%%%%%%%%%%%%%%%%%%%%%%%%%%%%%%%%%%%%%%%%%%%%%%%%%
% MACROS

    \newcommand{\lineunder} {
        \vspace*{-8pt} \\
        \hspace*{-18pt} \hrulefill \\
    }

    \newcommand{\header} [1] {
        \vspace{9pt}
        {\hspace*{-18pt}\vspace*{6pt} \large \textbf {#1}}
        \vspace*{-6pt} \lineunder
        \vspace{2pt}
    }

    \newcommand{\cvprofile}[7]{ %name,job,address,email,phone,blog,linkedin
        \begin{center}
            {\Huge \textbf {#1}}\\
            \vspace{5pt}
            {\large \textbf {#2}}\\
            \vspace{4pt}
            #3 $\cdot$\enspace #4 $\cdot$\enspace #5\\
            Blog: #6 $\cdot$\enspace LinkedIn: #7
            $\cdot$\enspace U.S. Citizen
%            LinkedIn: #7
        \end{center}
    }

    \newcommand{\employer}[4]{
        \vspace{3pt}
        \textbf{#1}  % company name
        \hfill #2\\  % city
        \textit{#3}  % title
        \hfill #4\\  % dates
        \vspace{3mm}
    }

    \newenvironment{jobtasklist}
        {
            \vspace{-12pt}
            \begin{itemize} \itemsep 0pt
        }{
            \end{itemize}
            \vspace{-3pt}
        }

    \newcommand{\university}[3]{
        \textbf{#1}  % name
        \hfill #2\\  % city
        #3\\         % degree
        \vspace{2mm}
    }

    \newenvironment{skillslist}
        {
            \hspace*{-0.07in}\begin{tabular}[t]{ l L{5.17in} }
        }{
            \end{tabular}
        }

    % underlines for important/emphasis
    \newcommand{\impt}[1]{\uline{#1}}


%%%%%%%%%%%%%%%%%%%%%%%%%%%%%%%%%%%%%%%%%%%%%%%%%%%%%%%%%%%%%%%%%%%%%
% VARIABLES

    % QR CODE (disabled unless QRCODE is enabled)
    %\usetag{QRCODE}

    % CV TYPE %%%%%%%%%%%%%%%%%%%%%%%%%%%%%%%%%%%%%%%%%%%%%%%%%%%%%%%
    \usetag{GISCV}

    % PERSONAL DETAILS
    %    CVname
    %    CVaddress
    %    CVemail
    %    CVblogurl
    %    CVlinkedin
    \input _private.tex

    % job details
    %   JobTitle
    %   JobCompany
    %   Hiring Manager
    %   Hiring Manager Title
    \input _company.tex

    % PDF METADATA
    \hypersetup{
        pdfauthor={\CVname},
        pdftitle={\CVname~- \JobTitle},
        pdfsubject={\JobTitle}
    }

%%%%%%%%%%%%%%%%%%%%%%%%%%%%%%%%%%%%%%%%%%%%%%%%%%%%%%%%%%%%%%%%%%%%%
% DOCUMENT

\begin{document}

%================== Profile ==================%
\tagged{QRCODE}{
    \AddToShipoutPicture*
        {\put(496,687)
            {\includegraphics[width=1.7cm,height=1.7cm]
            {linkedin-qr3.png}}}
    }

\vspace*{-30pt}

\cvprofile
    {\CVname}
    {\JobTitle}
    {\CVaddress}
    {\CVemail}
    {\CVphone}
    {\CVblogurl}
    {\CVlinkedin}

\vspace*{-6pt}



%================== Summary ==================%
\header{SUMMARY}
%Who , business functions / core competencies / unique value prop
\tagged{GISCV}{
{\JobTitle}
with 10+ years of experience performing complex GIS analysis and building web-based
geospatial applications for transportation, utility, and telecommunications infrastructure.
Skilled in Esri’s desktop and web GIS suite (ArcGIS Pro, ArcGIS Online, ArcGIS Enterprise),
Python, and FME Flow (Server), integrating geodatabases, web services,
and raster data through automated ETL pipelines on Azure.
Works directly with clients and subject matter experts to translate
business needs into data analysis requirements, delivering maps, dashboards,
and spatial REST APIs that support planning decisions.
}

%================== Skills ==================%
\header{SKILLS}
    \begin{taggedblock}{GISCV}
        \begin{skillslist}
            GIS Applications:&
                ArcGIS Pro, QGIS, GeoServer \linebreak
                ArcGIS Online, ArcGIS Enterprise, PostgreSQL/PostGIS \linebreak
            \\
            Programming \& Scripting:&
                Python, ArcPy, NumPy, Pandas, GeoPandas, Shapely, Rasterio, FastAPI \linebreak
                ArcGIS Maps SDK for JavaScript, React, HTML/CSS \linebreak
            \\
            ETL \& Data Architecture:&
                Python-based spatial ETL pipelines, spatial REST APIs \linebreak
                FME Flow (Server) \& FME Form, ArcGIS Pro ModelBuilder \linebreak
                Geodatabase management, spatial SQL functions \linebreak
            \\
            Cloud \& DevOps:&
                AWS (Certified Cloud Practitioner) \linebreak
                Azure (Certified Fundamentals) \linebreak
                Docker, containerization, Terraform IaC \linebreak
                GitHub, PyTest, Linux/bash scripting \linebreak
            \\
            Methodologies:&
                Agile, Jira sprint planning \linebreak
                Client needs analysis, cross-functional stakeholder engagement \linebreak

%            GIS Applications:&
%                ArcGIS, QGIS, ArcGIS Pro, ArcGIS Online, ArcGIS Enterprise \linebreak
%                - including Spatial Analyst, 3D Analyst, and Network Analyst extensions \linebreak
%                \\
%            Geodatabase management:&
%                PostgreSQL/PostGIS, Esri Enterprise, ArcGIS Online \linebreak
%               Proficient in \impt{Spatial modeling} to derive relationships \& patterns on the ground
%               Proficient in managing ArcGIS Online (AGOL) and developing with Experience Builder \linebreak
%               GeoServer administration, authorization, and configuration \linebreak
%               Proficient in Collecting, cleaning, and preprocessing data from multiple sources for analysis \linebreak
%               Proficient in processing geospatial data at scale\linebreak
%                \\
%            Spatial data analysis:&
%                %Extensive experience in creating, verifying \& converting between data types:
%                %    GeoJSON, Geodatabase, Geopackage \linebreak
%                Comprehensive experience in data collection, cleaning, consolidation and coordination \linebreak
%                Data Architecture, data modeling, warehousing and analytics \linebreak
%                \\
%            Programming \& Scripting:&
%                \impt{Python} libraries: ArcPy, NumPy, Pandas, GeoPandas, Shapely, Rasterio \linebreak
%                ArcGIS API for Python \linebreak
%                \impt{SQL}: spatial functions \linebreak
%                Spatial REST API development with FastAPI, Flask \linebreak
%                ArcGIS Pro Toolbox development with ArcPy \linebreak
%                Experience with \impt{machine learning} with PyTorch \linebreak
%                Python-based and FME spatial ETL pipelines \linebreak
%                JavaScript, .NET, HTML, CSS \linebreak
%                Orchestration with Airflow \linebreak
%                R programming and R Studio \linebreak
%                Spatial \impt{ETL Pipelines}: FME Flow \& FME Form \linebreak
%                Data Manipulation: Power Query, DAX \linebreak
%                Visualization \& Reporting: Microsoft Power BI \linebreak
%                \\
%            SQL:&
%                Spatial analysis \& geostatistics \linebreak
%                \\
%            Cloud Platforms:&
%                AWS Certified Cloud Practitioner \linebreak
%                Microsoft Certified: Azure Fundamentals \linebreak
%                Docker containerization and microservices \linebreak
%                Terraform IaC \linebreak
%                \\
%            Technical Proficiencies:&
%                Jira for task management, sprint planning, and issue tracking \linebreak
%                PostgreSQL - Spatial analysis, geostatistics, triggers, functions \linebreak
%                GitHub version control \linebreak
%                Unit testing with PyTest \linebreak
%                Linux and bash shell scripting \linebreak
%                Productivity \& Collaboration: Microsoft 365, Microsoft Teams, Jira \linebreak
%                Windows PowerShell \linebreak
%                Productivity \& Collaboration: Microsoft 365, Microsoft Teams \linebreak
%                \\
%            Communication \& collaboration: &
%                Experience in documenting problems related to GIS tools.
%                Extensive experience communicating technical information to technical and non-technical personnel
%                at various levels in the organization.
%                \\
%            ETL: &
%                Developing data pipelines and processes for cleaning, verifying,
%                and managing large data sets in Python
%                \\
%            Collaboration: &
%                Engaged with planners, clients, and partners to identify data
%                needs, plan data collection projects, refine analysis and models
%                \\
        \end{skillslist}
    \end{taggedblock}



%================== Experience ==================%
\header{WORK EXPERIENCE}



%------------------ Summit ------------------%
\employer
    {Summit Human Capital}
    {Remote}
    {GIS Data Engineer}
    {July 2025 -- Present}

\begin{jobtasklist}
    \item Translated utility client requirements into FME Flow (Server), Python ETL pipelines
            and REST APIs that process \impt{millions of infrastructure features daily}
            for an Esri Utility Network migration.
%    \item Partnered with utility stakeholders to analyze client needs
%            and optimize operational efficiency,
%            designing scalable FME/Python ETL pipelines and REST APIs
%            \impt{processing millions of daily features} during an Esri Utility Network migration.
    \item Developed custom ArcGIS Pro Python Toolboxes and automation scripts with
            the ArcGIS API for Python;
            added PyTest unit tests and GitHub-based CI/CD to speed releases and catch regressions early.
\end{jobtasklist}

%------------------ Picarro ------------------%
\employer
    {Picarro Inc.}
    {Santa Clara, CA}
    {GIS Data Engineer}
    {June 2024 -- April 2025}

\begin{jobtasklist}
    % Accomplished X as measured by Y by doing Z
%    \item Designed and implemented robust spatial ETL pipelines with Python, and ArcPy.
%            Integrated diverse geospatial datasets and to streamline data flow,
%            significantly \impt{enhancing data reliability} and reducing manual intervention.
    \item Automated CRS transformation, normalization, and validation of
            \impt{high-volume vector data} in Python before publishing it
            to ArcGIS Online web services via REST APIs.
%    \item Enhanced spatial data \impt{ETL pipeline management and efficiency},
%            improving data flow reliability and reducing manual intervention through custom Python,
%            ArcPy, and FME monitoring and recovery tools.
    \item Created a reusable Python library that standardized connections to corporate data sources and systems, 
            letting new ETL processes be built from shared components instead of custom code.
%    \item Supported agency-wide web mapping applications and data
%            consumption by \impt{optimizing data transfer processes} to ArcGIS Online
%            (AGOL) and ArcGIS Enterprise Portal, ensuring real-time data
%            availability for planning and public safety teams.
    \item Optimized geospatial \impt{data transfers of over 100,000 features daily}
            via FME Flow (Server), enabling real-time data 
            in enterprise web maps.

%    \item Optimized application APIs to enable \impt{high-volume geospatial data transfer} (hundreds of thousands of features daily)
%            from core databases to ArcGIS Online (AGOL) and ArcGIS Enterprise,
%            directly supporting dynamic web mapping applications and data consumption.
%    \item Enabled high-volume geospatial data transfer, successfully moving hundreds of thousands of
%            features daily from the core application database to ArcGIS Online (AGOL)
%            and ArcGIS Enterprise through the \impt{management and optimization of application APIs}.
\end{jobtasklist}

%\vspace{20pt}

\pagebreak

%------------------ Congruex ------------------%
\tagged{GISCV}{
    \employer
        {Congruex}
        {San Diego, CA}
        {GIS Analyst}
        {June 2020 -- June 2024}
    }

\begin{jobtasklist}
    % Accomplished X as measured by Y by doing Z
%    \item Collaborated with cross-departmental teams to gather requirements and deliver effective GIS solutions,
%            including supporting web-based GIS applications that
%            improved \impt{data accessibility for critical infrastructure design}.
    \item Led end-to-end deployment of a \impt{cloud GIS API platform} (Python/FastAPI on Azure),
            publishing 19 cataloged spatial REST web services;
            wrote end-user documentation and delivered demos and technical support to engineering teams.
%            delivering \impt{19 data analysis, reporting, and data management services}, providing critical business insights
%            and improving data accessibility for telecommunications infrastructure design.
%    \item Developed Python-based data loading tools (Jupyter Notebooks, GeoPandas, Rasterio, Folium) for the spatial data warehouse,
%            integrating \impt{automated quality assurance procedures} to ensure data accuracy and accelerate reliable data integration.
    \item Designed scalable Python-based tools for spatial data editing and management, using GitHub for collaborative code reviews,
            version control, and CI/CD best practices to ensure reliable production releases.
    \item Trained multiple technical and non-GIS users and \impt{led usability testing on web GIS applications};
            resolved GeoServer and PostgreSQL integration issues to keep spatial services available.
%    \item Developed and deployed Python-based tools and automations, including \impt{custom geoprocessing workflows} and data management scripts,
%            to enhance the speed and efficiency of spatial web application support and data integrity across GeoServer and PostgreSQL.
%    \item Improved the speed and efficiency of support staff tasks by streamlining
%            routine geoprocessing workflows; deployed Python-based tools and automations in
%            Jupyter Notebooks for \impt{data verification and management} across GeoServer and PostgreSQL.
%    \item Established reporting infrastructure
%            to provide geospatial reports and analyses on
%            scheduled basis to support project management and administrative business processes,
%            deployed on Azure with Docker and Python
%    \item Developed reusable library in Python to automate the configuration and management of
%            corporate mapping platform comprising more than 7000 lines of code
    \item Engineered an automated, containerized spatial ETL pipeline on Azure
            (Python, GeoPandas, Fiona, Rasterio, Docker) incorporating custom spatial intersect
            validation and QA algorithms to resolve boundary anomalies
            across \impt{13 million features per day}.
%    \item Developed and implemented an automated, containerized spatial ETL pipeline on Azure
%            for \impt{data cleaning and quality assurance},
%            consolidating geospatial data from partner organizations,
%            processing over 13 million features daily, and ensuring data integrity.
%    \item Facilitated \impt{end-user usability testing} with both technical and non-GIS stakeholders,
%        effectively translating complex geospatial concepts into actionable feedback for platform improvements and acting
%        as a critical bridge between departments.
%    \item Managed ETL pipeline and API documentation in Confluence along with support requests in Jira
\end{jobtasklist}

%------------------ Freelance ------------------%
\tagged{GISCV}{
    \employer
        {Freelance}
        {San Diego, CA}
        {GIS Data Analyst}
        {January 2019 -- May 2020}
    }

\begin{jobtasklist}
%       Accomplished X as measured by Y by doing Z
%        \item Developed and deployed targeted surveys, providing insights into
%                customer engagement and product utilization.
        \item Worked directly with clients to translate business needs
                into GIS data analysis requirements,
                delivering maps and survey analytics that improved campaign effectiveness.
        \item Built \impt{interactive web GIS dashboards}
                (JavaScript, HTML/CSS, Web APIs)
                to visualize survey results for executive clients,
                while automating end-to-end data cleaning pipelines using Python
                to reduce data errors.
\end{jobtasklist}


%------------------ BC Transit ------------------%
\employer
    {BC Transit}
    {Victoria, Canada}
    {GIS Analyst}
    {January 2013 -- October 2018}

\begin{jobtasklist}
    \item Collected and edited spatial data for \impt{urban planning analyses};
            produced digital and hard-copy maps and capital-improvement
            recommendations for executives and regional stakeholders.
    \item Implemented a Python-based spatial ETL pipeline
            \impt{processing over 150,000 transactions per day}
            from a fleet of 700+ vehicles, providing detailed insights into transit ridership patterns.
%    \item Developed a \impt{performance metrics system} to monitor the security and
%            quality of regional reporting systems, improving alerting and
%            notification for farebox security incidents and technical problems.
%    \item Created a continuous monitoring system with multiple departments,
%            ensuring an uninterrupted stream of ridership data that provided
%            up-to-date statistical reports and analysis.
%    \item Conducted advanced data modeling and analysis on complex ridership data, utilizing sophisticated
%            GIS algorithms to generate insightful reports that directly supported strategic decision-making and
%            demonstrated the ability to \impt{translate raw data into actionable intelligence} for stakeholders solving real-world transportation
%            problems.
    \item Performed \impt{complex GIS analysis of transportation and ridership data} to
            inform policy and guide regional transit service planning.
%        \item Conducted ongoing GIS data maintenance, manipulation, analysis, extraction, and generation
%                and managing organizational geospatial data warehouse.
    \item Received BC Transit’s 2015 Recognizing \impt{Excellence and Values Award}
            for contributing to a cross-departmental team that significantly
            improved transit vehicle farebox data collection.
\end{jobtasklist}


\vspace{25pt}

%================== Education ==================%
\header{EDUCATION}
    \university
        {Queen's University}
        {Kingston, Canada}
        {Master of Science, Computer Science}

    \university
        {University of Ottawa}
        {Ottawa, Canada}
        {Graduate Diploma, International Development}

    \university
        {Queen's University}
        {Kingston, Canada}
        {Bachelor of Science, Psychology \& Computer Science}
    \university
        {Udacity Nanodegree (Bertelsmann Scholarship)}
        {}
        {AI Programming with Python}
    \university
        {Udacity Nanodegree (Bertelsmann Scholarship)}
        {}
        {Cloud Architect using Microsoft Azure}

\end{document}
